We extend the static decision quotient framework of Paper 4 to stochastic and sequential regimes. Given a decision problem with actions $A$, factored state space $S = X_1 \times \cdots \times X_n$, and utility $U$, we ask: what complexity is required to identify sufficient coordinate sets when (a) the state is drawn from a distribution $P$, or (b) decisions unfold over time with transitions and observations?

\textbf{Results:}
\begin{itemize}
\item \textbf{Stochastic regime}: STOCHASTIC-SUFFICIENCY-CHECK is PP-complete via MAJSAT reduction. Tractable under product distributions, bounded support, log-concave.

\item \textbf{Sequential regime}: SEQUENTIAL-SUFFICIENCY-CHECK is PSPACE-complete via TQBF reduction. Tractable under full observability, bounded horizon, tree structure.

\item \textbf{Regime hierarchy}: Static $\subset$ Stochastic $\subset$ Sequential; coNP $\subset$ PP $\subset$ PSPACE.

\item \textbf{Substrate cost}: The integrity-competence verdict is substrate-independent; trajectories diverge based on substrate cost $\kappa$.

\item \textbf{Temporal learning}: Bayesian structure detection reduces abstention over time.

\item \textbf{Temporal integrity}: Evidence-monotone retractions preserve integrity across sequences.
\end{itemize}

The reduction constructions are machine-checked in Lean 4 (\LeanTotalLines\ lines, \LeanTotalTheorems\ theorems).

\textbf{Keywords:} computational complexity, decision theory, stochastic decision problems, POMDPs, polynomial hierarchy, PSPACE
